% =============================================================================
% CONCEPTUAL ITEM: what representation independence actually OBLIGES you to
% prove. Students can usually recite "the client can't tell" -- this asks for
% the shape of the argument instead: a relation, preserved by the builder, under
% which observations agree.
%
% Deliberately NOT the same punchline as mod2.tex (which asserts that two queue
% implementations are indistinguishable). Here the answer is the proof
% obligation, not the conclusion.
%
% Three operations only -- one seed, one builder, one observation -- which is
% the minimum that still forces all three clauses. Using ACC rather than SET
% also makes the two representations genuinely different TYPES (int vs
% int list), so "relation, not bijection" is immediate: every list summing to 7
% relates to the single accumulator 7.
%
% NOTE the asymmetry that students get wrong: clauses (1) and (2) are about
% R-relatedness, but clause (3) is about EQUALITY. `total` returns int, not the
% abstract t, so R does not even typecheck there.
%
% Both implementations verified to agree in SML/NJ (probe: add 10, ~7, 4).
% =============================================================================

\part[4]\hfill

  Two structures \code{A1} and \code{A2} implement the same signature
  \code{ACC}.

  \begin{codeblock}
    signature ACC =
      sig
        type t
        val init  : t
        val add   : int -> t -> t
        val total : t -> int
      end
  \end{codeblock}

  We want to show no client written against \code{ACC} can distinguish
  \code{A1} from \code{A2} --- in other words, they are equivalent through
  representation independence.

  Given a relation $R(\code{a1}, \code{a2})$ relating values \code{a1 : A1.t}
  and \code{a2 : A2.t}, write down the facts you would need to prove. Be
  careful to state each one as either an $R$-relatedness or an equivalence,
  whichever is appropriate.

  \begin{solutionorbox}[22em]
    Checking example clients cannot work --- there are infinitely many. Instead
    we show $R$ holds initially and is preserved, and that observations agree:

    \begin{enumerate}
      \item \textbf{The seeds are related:} $R(\code{A1.init}, \code{A2.init})$.

      \item \textbf{\code{add} preserves $R$.} If $R(\code{a1}, \code{a2})$,
      then for every value \code{x : int},
      $$R(\code{A1.add x a1},\; \code{A2.add x a2})$$

      \item \textbf{Observations are equal on related states.} If
      $R(\code{a1}, \code{a2})$, then
      $$\code{A1.total a1} \;\eeq\; \code{A2.total a2}$$
    \end{enumerate}

    Note the asymmetry between (2) and (3). \code{add} returns the
    \textbf{abstract} type \code{t}, so the best we can say is that its results
    are \textit{related}. But \code{total} returns \code{int}, a type the client
    can see and compare directly --- so its results must be genuinely
    \textbf{equivalent}. (Writing $R$ in (3) would not even typecheck: $R$
    relates a \code{A1.t} to a \code{A2.t}, not two \code{int}s.)

    Together these give an induction on the client's sequence of calls: the two
    structures start related, every \code{add} keeps them related, and any
    \code{total} the client performs returns the same answer either way. Since
    \code{t} is abstract, building with \code{init}/\code{add} and observing
    with \code{total} is all a client can do.
  \end{solutionorbox}

\newpage
